%% appendix_dragonfly.tex — Dragonfly Mission Validation
%% Moved from Discussion chapter to appendix for word-count compliance.

\chapter{Dragonfly Mission: Site Validation and Instrument Mapping}
\label{app:dragonfly}

\section{Mission Science Objectives}
\label{sec:dragonfly_objectives}

NASA's \dragonfly{} rotorcraft \citep{Lorenz2018Dragonfly}, due at \titan{} in 2034, traverses $\sim$175\,km from Belet dunes to Selk crater.  Its instruments --- DraMS (mass spectrometry), DraGNS (gamma-ray and neutron spectrometry), DragonCam (imaging), and DraGMet (meteorology) --- directly measure the features underpinning the habitability scores at its landing sites.

\section{Two Distinct Habitability Frameworks}
\label{sec:two_frameworks}

Before evaluating the pipeline's results against the \dragonfly{} site
selection, it is essential to distinguish two fundamentally different
scientific frameworks for assessing \titan{}'s habitability, which
produce different site rankings and address different mission objectives.

\paragraph{Present-day solvent habitability.}
This framework asks: \textit{where does liquid solvent exist at the present
epoch that could support metabolic chemistry?}  For water-based life, the
answer is the inaccessible subsurface ocean.  For non-aqueous hydrocarbon
biochemistry \citep{McKaySmith2005}, the answer is the north polar
methane--ethane lakes.  The Bayesian pipeline in this thesis implements
\textit{present-day solvent habitability}: the eight features weight liquid
hydrocarbon availability ($w_1 = 0.25$), methane cycle activity
($w_4 = 0.15$), and surface--atmosphere interaction ($w_5 = 0.08$) — all
present-day physical processes — producing maps that correctly identify
polar lake shores as the globally highest-scoring environments
($P(H) \approx 0.57$ at the present epoch; Freeman Lacus).

\paragraph{Prebiotic chemistry habitability.}
This framework asks: \textit{where has liquid water most plausibly
contacted organic material on geologically relevant timescales?}
\citet{Neish2018} identified large impact craters as the clear answer.
The thermal energy of a major impact melts the water-ice crust, producing
transient liquid water that remains liquid for $10^3$--$10^6$\,yr at
temperatures that accelerate chemical reaction rates by several orders of
magnitude over ambient conditions \citep{Artemieva2003, OBrien2005}.
Water-ice mixed with \titan{}'s photolysed tholin organics under these
conditions is expected to produce amino acids and nucleobase precursors
via Strecker synthesis \citep{Neish2010}.  Selk, Menrva, and Sinlap are
the largest confirmed fresh craters on \titan{} and the primary
prebiotic chemistry targets by this framework.  This is the scientific
framework within which \dragonfly{} was selected: the mission explicitly
targets \textit{evidence of past water--organic contact chemistry}, not
present-day liquid \citep{Lorenz2021Landing}.

\paragraph{Why the two frameworks give different site rankings.}
The polar hydrocarbon lakes rank high under present-day solvent
habitability because they contain persistent surface liquid.  They rank
\textit{lower} under the prebiotic chemistry framework precisely because
liquid methane and ethane are non-polar: they do not dissolve the ionic
species required for Strecker synthesis, and they never contact the
water-ice crust in a chemically productive way.
Conversely, impact craters rank low under present-day solvent habitability
(no present liquid at the equatorial surface) but high under prebiotic
chemistry habitability (confirmed past liquid water).  \citet{Neish2018}
explicitly concluded that the polar lakes are \textit{not} the best places
to search for biosignatures, because the non-polar solvent cannot drive the
relevant prebiotic chemistry.

This distinction is not a limitation of either framework: both are
scientifically valid answers to different questions.  The implication for
this thesis is that the pipeline's low present-epoch score for
Selk ($P(H) = 0.243$) is not a contradiction of the \dragonfly{} mission's rationale
--- it is a \textit{consequence} of targeting different habitability.
Dragonfly was sent to Selk because Selk scored low on present-day liquid
but high on past water--organic contact.

\section{Quantitative Validation of Site Selection}
\label{sec:dragonfly_quant}

The pipeline is a present-day solvent habitability model.  By this
metric, the north polar lake margins are the most habitable environments
on \titan{}, reaching $P(H) \approx 0.57$ (Freeman Lacus) and
$0.53$ (Oib Lacus).  Selk scores $P(H) = 0.243$, below the
prior mean of $0.33$, because it lacks present-day surface liquid
($f_1 \approx 0.05$) and its methane-cycle and surface--atmosphere scores
are equatorial-low.  \textit{This is scientifically correct}: a model that
correctly captures present-day solvent habitability should rank polar lake
shores above a dry equatorial crater.

The independent validation comes from the geomorphological diversity
feature.  Selk's $f_7 = 0.659$ is the highest value of any equatorial
site and 7.2$\times$ the global median of 0.092 --- a result derived from
Cassini SAR terrain data alone, without reference to mission planning.
Geomorphological diversity captures the simultaneous presence of four
terrain classes (dune, plains, labyrinth, crater floor) within 45\,km,
which is the Cassini-era observational proxy for the organic--water-ice
mixing interface that defines the prebiotic chemistry target.  The
pipeline independently identifies Selk as the equatorial site with the
highest terrain-class heterogeneity --- precisely the physical attribute
cited in the \dragonfly{} science rationale
\citep{Lorenz2021Landing, Neish2018}.

The polar lake margins represent a complementary, not competing, science
target.  A future mission targeting present-day habitability --- such as
the POSEIDON lake lander concept \citep{Tosi2022POSEIDON} --- would be
directly guided by the polar-lake-dominant result of this pipeline.

\section{Instrument--Feature Correspondence}
\label{sec:instruments}

DraMS atmospheric and surface composition measurements constrain $f_3$ (acetylene abundance) through depletion relative to photochemical model predictions, and $f_2$ (organic inventory) via amino acid and HCN polymer detection.  DraGNS measures bulk hydrogen abundance, constraining $f_8$ (subsurface water proximity).  The DraGMet EFIELD experiment directly measures the surface electric field, constraining $f_4$.

Enantiomeric excess in amino acids (expected $>$10\% if biogenic vs.\ racemic for abiotic Strecker products) is the highest-priority biosignature for DraMS.  Acetylene depletion below the photochemical steady-state prediction \citep{McKaySmith2005} would constitute strong but non-unique evidence for acetylenotrophy.  Carbon isotopic fractionation ($\delta^{13}$C depletion of $\sim$30--50\,\textperthousand{} relative to atmospheric \ce{CH4}) would be consistent with biological methane cycling.  All three biosignatures have plausible abiotic explanations and must be interpreted jointly.

\section{Interpreting a Null Result}
\label{sec:null_result}

A null biosignature result at Selk would constrain specifically the
``preserved impact-melt Strecker chemistry'' pathway at that location and
epoch.  It would not constrain: (1) acetylenotrophy at the north polar lake
margins, which score $\sim$30 percentage points higher than Selk in the present
epoch and which \dragonfly{} does not visit; (2) life in the subsurface
ocean, insensitive to surface sampling; or (3) extinct life from the LHB
whose signatures may be buried below accessible depth.  The habitability maps
provide a rigorous framework for reporting a null result in terms of which
specific pathways have been tested and which remain open.
